\documentclass[11pt,a4paper]{article}

\usepackage[utf8]{inputenc}

\usepackage[T1]{fontenc}

\usepackage[polish]{babel}

\usepackage{lmodern}

\usepackage{geometry}

\geometry{margin=27mm}

\usepackage[hidelinks]{hyperref}

\title{Kalibracja polegania na wskazówkach AI w ekranie 2D, XR i symulowanym humanoidzie}

\author{Mateusz Pomianek}

\date{Szkic naukowy, wersja 1.0.0 | 28 listopada 2026}

\begin{document}

\maketitle

\begin{abstract}

Proponowana miara główna to interakcja medium i poprawności porady dla decyzji o jej przyjęciu.

\end{abstract}

\section{Wprowadzenie}

Zaufanie może wpływać na decyzję o poleganiu na automatyzacji, a błędne poleganie i niewykorzystanie trafnej pomocy mają różne konsekwencje (Lee \& See, 2004; Parasuraman \& Riley, 1997).

\section{Cel i pytania badawcze}

Hipotezy robocze wymagają zatwierdzenia przed rozpoczęciem właściwej weryfikacji.

\paragraph{H1.} Różnica między prawdopodobieństwami przyjęcia poprawnej i błędnej porady będzie zależna od medium przy identycznej treści i jakości rekomendacji.

\paragraph{H2.} Symulowany humanoid zwiększy przyjmowanie błędnych porad względem interfejsu 2D przy identycznych informacjach o systemie.

\section{Materiał i metody}

Pierwszorzędowy wskaźnik: interakcja medium i poprawności porady dla decyzji o jej przyjęciu.

Główny test obejmie interakcję medium i poprawności w modelu logistycznym z efektami uczestnika i scenariusza. Kolejność mediów wymaga zrównoważenia, a samą treść, informacje o trafności i dostępność przesłanek należy utrzymać na tym samym poziomie. Wynik powinien oddzielać deklarowane zaufanie od decyzji.

Przed badaniem z udziałem ludzi należy zatwierdzić wymagania etyczne, zasady zgody i ochronę danych.

\section{Plan analizy}

Zamrożenie estymandu, reguł oceny jakości i obsługi braków nastąpi przed dostępem do danych wynikowych. Raport obejmie oszacowanie kontrastu głównego z przedziałem niepewności i pełne mianowniki. Jeśli przewidziano hipotezy drugorzędowe, przed analizą należy ustalić sposób kontroli wielokrotności.

\section{Wyniki}

W niniejszym pliku nie umieszczono wyników empirycznych, statystyk ani nieprzeprowadzonych analiz. Sekcja wymaga uzupełnienia po uzyskaniu i weryfikacji materiału dowodowego.

\section{Dyskusja i ograniczenia}

Ograniczenia obejmują możliwe błędy pomiaru, dobór próby, alternatywne mechanizmy i ograniczoną możliwość uogólniania. Korelacja nie dowodzi przyczynowości; poprawność techniczna nie jest dowodem znaczenia klinicznego.

\section{Wnioski}

Hipotezy i sposób pomiaru stanowią podstawę oceny wykonalności oraz przygotowania właściwego protokołu.

\section*{Bibliografia}

{\setlength{\parindent}{0pt}
\setlength{\parskip}{5pt}
\hangindent=1.27cm
Lee, J. D., \& See, K. A. (2004). Trust in automation: Designing for appropriate reliance. \textit{Human Factors, 46}(1), 50--80. \url{https://doi.org/10.1518/hfes.46.1.50_30392}\par
\hangindent=1.27cm
Parasuraman, R., \& Riley, V. (1997). Humans and automation: Use, misuse, disuse, abuse. \textit{Human Factors, 39}(2), 230--253. \url{https://doi.org/10.1518/001872097778543886}\par
}

\end{document}
